\documentclass[a5paper,10pt]{article}

% https://github.com/InDevRus/filippov-solutions

% Нумерация страниц
\pagenumbering{gobble}

% Настройка полей
\usepackage{geometry}
\geometry{left=1cm}
\geometry{right=1cm}
\geometry{top=1cm}
\geometry{bottom=1.5cm}

% Вставка таблицы
\usepackage{array}
\usepackage{tabularx}

% Инструменты выравнивания формул
\usepackage{amsmath}
\usepackage{mathtools}

% Диакритика в математике
\usepackage{amsfonts}

% Вычисления размеров на ходу
\usepackage{calc}

% Удобные записи производных
\usepackage[italic=true]{derivative}

% Настройки сносок
\usepackage{enotez}

% Длинные знаки векторов
\usepackage{anyfontsize}
\usepackage[b]{esvect}

% Вставка картинок
\usepackage{graphicx}

% Вставка красной строки
\usepackage{indentfirst}
\setlength{\parindent}{1em}

% Условия в командах
\usepackage{xifthen}

% Луа-скрипты
\usepackage{luacode}

% Настройка команд отдельно в математическом и текстовом режимах
\usepackage{mathcommand}

% Для повторения LaTeX кода
\usepackage{multido}

% Конфигурация отступов формул
\delimitershortfall-1sp
\usepackage{mleftright}
\mleftright

% Растягивание объектов
\usepackage{relsize}

% Обтекание текстом
\usepackage{wrapfig}
\usepackage{subcaption}
\usepackage{floatflt}

% Поддержка ввода кириллицы
\usepackage[english, russian]{babel}

% Настройка корневой позиции для компилятора
\usepackage{import}
\providecommand*{\ProjectRootPath}{..}

\import{\ProjectRootPath/preambles/macros/}{theorems}

\import{\ProjectRootPath/preambles/fonts}{font_setup}

% Знак градуса
\usepackage{siunitx}

\import{\ProjectRootPath/preambles/macros/}{operators}
\import{\ProjectRootPath/preambles/macros/}{redefinitions}

\import{\ProjectRootPath/preambles/fonts}{letters_setup}
\import{\ProjectRootPath/preambles/fonts/adjustments}{custom_superscripts}

\import{\ProjectRootPath/preambles/custom}{formatting}
\import{\ProjectRootPath/preambles/custom}{frames}
\import{\ProjectRootPath/preambles/custom}{list_setup}

% TODO: перевести команды на современный стиль объявления

\begin{document}
    $x \parder{z} {x} + 2y \parder{z} {y} = \pow{x}{2} y + z$.
    $\dfrac {dx} {x} = \dfrac {dy} {2y} = \dfrac {dz} {\pow{x}{2} y + z}$.

    $\dfrac {dx} {x} = \dfrac {dy} {2y}$
    представляет собой интегрируемую комбинацию.
    $2\ln\vbr{x} = \ln\vbr{y} + A$.
    $\dfrac {y} {\pow{x}{2}} = A$.
    Далее подставим в
    $\dfrac {dx} {x} = \dfrac {dz} {\pow{x}{2} y + z}$ первый интеграл $y = A\pow{x}{2}$:
    $\dfrac {dz} {dx} = \dfrac {\pow{x}{2} y + z} {x} = xy + \dfrac {z} {x} = \linebreak = A\pow{x}{3} + \dfrac {z} {x}$.
    Получилось линейное уравнение первого порядка.
    $xz\' - z = A\pow{x}{4}$.
    $\dfrac {xz\' - z} {\pow{x}{2}} = \linebreak = \dfrac {d} {dx} \br{\dfrac {z} {x}} = A\pow{x}{2}$.
    $\dfrac {z} {x} = \dfrac {A} {3} \pow{x}{3} + B$.
    $B = \dfrac {z} {x} - \dfrac {A} {3} \pow{x}{3} = \dfrac {z} {x} - \dfrac {xy} {3}$.

    Итак, общее решение имеет вид
    $F\br{\dfrac {y} {\pow{x}{2}}; \dfrac {z} {x} - \dfrac {xy} {3}} = 0$.
    $\dfrac {z} {x} - \dfrac {xy} {3} = f\br{\dfrac {y} {\pow{x}{2}}}$. В явном виде искомые функции имеют вид
    $z\br{x; y} = \dfrac {\pow{x}{2} y} {3} + x f\br{\dfrac {y} {\pow{x}{2}}}$.
\end{document}
